\section{Mixed-parity alpha constructions}
\label{mc:chapter}
\label{app:ten-alpha}

This section supplies the mixed-parity constructions used in the
seven-arm arithmetic cover and the even-minimum eight-arm theorem.
All arm lengths count edges, and all displayed words omit their common
initial vertex, whose label is zero.  Intervals in formulas for labels
or colors are integer intervals.  The prescribed-endpoint path inputs
are exactly Lemmas~\ref{common:even-path} and~\ref{common:odd-path}.
Their existence is classical; all subsequent operations and inventories
are proved here.  In particular, an arbitrary alpha labeling is never
assumed to have a prescribed Walecki arm.

\subsection{A shortening operation with a prescribed arm}
\label{mc:shortening}

\begin{lemma}[Maximum-Walecki shortening]\label{mc:shorten}
Suppose a root-zero alpha labeling of a spider with $N$ edges contains
the arm
\[
 W=(N,1,N-1,2,\ldots,N-s)
\]
of length $2s+1$, and every label on the other arms is below $N-s$.
Deleting the last leaf of $W$ and subtracting one from every remaining
label above $N-s$ gives a root-zero alpha labeling with that arm
shortened to $2s$.  The alpha index is unchanged.
\end{lemma}
\begin{proof}
Write $h=N-s$.  Label compression is a bijection onto
$\interval{0}{N-1}$.  Only labels on $W$ are changed.  The deleted
edge has color $N-2s$; the remaining colors of $W$, formerly
$\interval{N-2s+1}{N}$, each decrease by one.  They now give
$\interval{N-2s}{N-1}$, adjoining the unchanged colors
$\interval{1}{N-2s-1}$ on the other arms.  Every surviving upper
label on $W$ is at least $h$, still above the old cut.  Lower labels
and the root are unchanged, so every edge still crosses that cut.
\end{proof}

\subsection{Odd--even--odd triples and their shortening}
\label{mc:oeo}

\begin{proposition}\label{mc:oeo-central}
Let $a<b<c$ have respective parities odd, even, odd.  If
$c\le2a-3$, then $\Sp{a,b,c}$ has a root-zero alpha labeling
with its shortest arm equal to the maximum Walecki word.
\end{proposition}
\begin{proof}
Write $(a,b,c)=(2s+1,2t,2u+1)$, so
$1\le s<t\le u\le2s-1$.  Put
\[
 S=u-t,\qquad T=u-s,\qquad P=s+t-u,\qquad v=P+1.
\]
Thus $0\le S<T$, $P\ge T-S+1$, and $1\le v\le u$.
Take an even-order alpha path on $\interval{1}{2T}$ from upper
endpoint $2T-S$ to lower endpoint $T-S$.  Raise its uppers by one.
It now omits label $T+1$ and has colors $\interval{2}{2T}$.
Raise its uppers by a further $2P$.  Before it place an even-order
alpha path on $\interval{1}{2P}$ from $P+T-S+1$ to $T-S+1$,
translated by $T$.  Both requests are legal since
$1\le T-S+1\le P$.  The join has color $2P$.

The resulting path $Q$ has lower labels $\interval{1}{t}$,
upper labels $\interval{t+1}{2t+1}\setminus\{t+v\}$, and colors
\[
 \interval{1}{2t}\setminus\{2P+1\}.
\]
Its first and last labels are $E=2t-s+1$ and $F=t-s$.
These assertions follow from the adjacent prefix intervals: the
prefix contributes colors $\interval{1}{2P-1}$ and the old path
contributes $\interval{2P+2}{2t}$, besides their joining color.

Map the lowers of $Q$ by $x\mapsto x+s$ and the uppers by
$x\mapsto x+s+2u$.  Call this arm $A$.  Its internal colors are
$\interval{2u+1}{2u+2t}\setminus\{2u+2P+1\}$ and its root
color is $2u+2t+1$.
For arm $B$, take an even-order alpha path on
$\interval{1}{2u}$ from $u+v$ to $v$, append $2u+v$, and
translate every label by $s+t$.  Its internal colors are
$\interval{1}{2u}$, and its root color is
$s+t+u+v=2u+2P+1$.

The lower labels of $A$ and $B$ are respectively
$\interval{s+1}{s+t}$ and $\interval{s+t+1}{s+t+u}$.
The upper block of $B$ is $\interval{s+t+u+1}{s+t+2u}$;
that of $A$ is $\interval{s+t+2u+1}{s+2t+2u+1}$ with the
single hole filled by the appended leaf of $B$.
Set $N=2s+2t+2u+2$ and adjoin
$W=(N,1,N-1,2,\ldots,N-s)$.  Its labels fill
$\interval{1}{s}\cup\interval{N-s}{N}$, and its colors are
$\interval{2t+2u+2}{N}$.  All labels and colors are now consecutive
without repetition.  All edges cross $s+t+u$, and every label off
$W$ is below its last leaf.  This proves the pinning as well.
\end{proof}

\begin{corollary}\label{mc:eeo-central}
If $a,b$ are even, $6\le a<b<c$, $c$ is odd, and
$c\le2a-1$, then $\Sp{a,b,c}$ is root-zero alpha.
\end{corollary}
\begin{proof}
Apply Proposition~\ref{mc:oeo-central} to $(a+1,b,c)$:
the parity inequalities imply $a+1<b<c$, and
$c\le2(a+1)-3$.  Its explicitly prescribed shortest arm can be
shortened by Lemma~\ref{mc:shorten}.  This proves exactly the bound
$c\le2a-1$; a larger bound is not inferred from this operation.
\end{proof}

\subsection{Even--odd--odd central and repeated triples}
\label{mc:eoo-central-section}

We use the upper-hole kernels $K(S,U,T,V)$ and their completion from
Definition~\ref{ok:kernel-definition} and Lemma~\ref{ok:completion}.
For clarity, their final completion has labels
\[
 \interval{1}{T}\ \cup\
 \bigl(\interval{T+1}{2T+2}\setminus\{T+V\}\bigr),
\]
missing color $S+T-U+V$, and ordered endpoints
$2T-S+2$, $2T+2S-U-V+3$.  Its three completed arm lengths are
$(2S+1,2T+1,2U+1)$.  That operation explicitly supplies the
maximum-Walecki shortest arm.  The auxiliary source kernels used in
growth need not themselves satisfy the bounds required for a final
completion.

\begin{lemma}[Pinned central AP kernels]\label{mc:eoo-ap}
For $d\ge1$, $s\ge2d+2$, the spider
\[
 \Sp{2s,2s+2d+1,2s+4d+1}
\]
is root-zero alpha.  For $d=1$ the stronger range $s\ge2$ holds.
\end{lemma}
\begin{proof}
For $d\ge2$, the following explicit word is a kernel
$K(d,3d,2d,2)$:
\[
\begin{split}
 &(3d+2),\\
 &(d+1-j,3d+3+j)\quad(0\le j\le d-1),\\
 &(1,2d+1),\\
 &(2d+1-j,2d+2+j)\quad(1\le j\le d-1).
\end{split}
\]
Its lowers partition $\interval{1}{2d}$; its uppers partition
$\interval{2d+1}{4d+2}\setminus\{2d+2\}$.  The endpoints
are $3d+2,3d+1$.  The word through its first label $1$ gives
colors $\interval{2d+1}{4d+1}$, the next edge gives $2d$, and
the final block, including its join, gives
$\interval{1}{2d-1}\setminus\{2\}$.  Thus its missing color
is precisely $2$.  The restriction $d\ge2$ excludes an endpoint
equal to the label hole.

Apply diagonal growth (Lemma~\ref{ok:diagonal-growth}) with
$P=s-d\ge d+2$.  In explicit terms its even-order path has upper
endpoint $P+d+2$ and lower endpoint $d+2$, is translated by $2d$,
and is put before the kernel after raising the latter's uppers by
$2P$.  The resulting parameters are
\[
 (S,U,T,V)=(s,s+2d,s+d,s-d+2).
\]
Here $1\le V\le U$.  Complete this kernel, then shorten its
prescribed Walecki arm by Lemma~\ref{mc:shorten}.  The resulting
index is $3s+3d$ and the arm lengths are those claimed.

For $d=1$, use these three literal kernels, in each case with
$T=s+1$ and $U=s+2$:
\begin{center}\begin{tabular}{ccp{0.72\textwidth}}
$s$&$V$&word\\\midrule
2&2&$(6,1,8,2,4,3,7)$\\
3&2&$(7,4,5,3,8,2,9,1,10)$\\
4&4&$(8,2,12,1,10,5,6,4,7,3,11)$.
\end{tabular}\end{center}
Direct subtraction checks the full kernel inventories.  For every
$s\ge5$, grow the first kernel by $P=s-2\ge3$, using even-path
endpoints $P+3,3$.  This gives parameters $(s,s+2,s+1,s)$.
Completion and shortening apply in all cases.
\end{proof}

\begin{proposition}\label{mc:eoo-central}
If $a\ge6$ is even, $b,c$ are odd, and
$a<b\le c\le2a+1$, then $\Sp{a,b,c}$ is root-zero alpha.
\end{proposition}
\begin{proof}
First assume $c\le2a-3$.  The odd source $(a+1,b,c)$ lies
in the pinned odd central construction of
Lemma~\ref{ok:central-kernel}, apart from its AP plane.
At $a+1=b=c$ use the explicit equal-odd construction in
Proposition~\ref{ok:equal-odd-three}; at $a+1<b=c$ one may use
the pinned repeated kernels of Lemma~\ref{ok:repeated-kernels}.
On the strict AP plane $a+1+c=2b$, write the lengths as
$(2s,2s+2d+1,2s+4d+1)$; its central inequality is
$s\ge2d+2$, so Lemma~\ref{mc:eoo-ap} supplies the missing pin.
Shortening the stated arm proves this first range.

For the two remaining odd layers write $a=2s$, $b=2t+1$,
$s\ge3$.  At $c=4s-1$ the source largest half is $2s-1$.
If $t\le2s-3$, use the explicit odd-prefix boundary construction
of Lemma~\ref{ok:odd-prefix-boundaries}.  If $t=2s-2$,
$s\ge4$, use Lemma~\ref{ok:gap-two-kernels}.  The exceptional
source $(7,9,11)$ uses the $s=3,d=1$ literal kernel of
Lemma~\ref{mc:eoo-ap}, rather than an unpinned AP existence
statement.  At $t=2s-1$, the explicit last-ray repeated kernel
has $V=s$; its even completion endpoint is $3s-1$ on
$\interval{1}{4s-2}$, an allowed request.

At $c=4s+1$ the reversed boundary construction of
Lemma~\ref{ok:reverse-boundary} covers $s\le t\le2s$.
Its repeated endpoint $t=2s$ is the explicit twice-ray kernel with
$V=s+3$ and even completion endpoint $3s+3$ on
$\interval{1}{4s}$; this is allowed since $s\ge3$.
Every cited construction completes by actually adjoining the standard
maximum-Walecki arm.  Apply Lemma~\ref{mc:shorten} in each case.
\end{proof}

\begin{proposition}\label{mc:eoo-repeated}
For every even $a\ge6$ and odd $b>a$, $\Sp{a,b,b}$ is
root-zero alpha.
\end{proposition}
\begin{proof}
Write $a=2s,b=2t+1$.  The pinned all-odd repeated source has
$V=s+1$ when $s<t\le2s-2$ or $t\ge2s+1$, $V=s$ at
$t=2s-1$, and $V=s+3$ at $t=2s$.  At $t=s$ use the explicit
equal-odd construction with its first Walecki arm.  These are the
complete repeated constructions of Proposition~\ref{ok:equal-odd-three}
and Lemma~\ref{ok:repeated-kernels}, and exhaust $t\ge s$.
All completion requests $t+V$ lie in $\interval{t+1}{2t}$ in
their stated domains.  Shortening the pinned arm gives the result.
In particular $b=a+1$ is included without a separate inference
about pinning of arbitrary all-odd alpha labelings.
\end{proof}

\subsection{Equal-even near triples}
\label{mc:eeo-near-section}

\begin{proposition}\label{mc:eeo-near}
For every even $a\ge6$, $\Sp{a,a,a+1}$ is root-zero alpha.
\end{proposition}
\begin{proof}
There is a useful endpoint operation.  In an $N$-edge root-zero
alpha labeling, suppose an even arm has first label $H$, last label
$l$, and $H+l=N+1$.  Replace its word $P$ by
$(N+1,\operatorname{rev}(P))$.  Reversal retains all old internal
colors.  The old root color $H$ is replaced by the joining color
$(N+1)-l=H$, and the new root supplies color $N+1$.
Alternation and the cut are retained, since the reversed even word
starts lower and ends upper.  This proves the operation.

Write $a=2r$.  If $r\ge7$ is odd, the appropriate equal-even construction in
Section~\ref{templates:equal-even} has complementary endpoints
$H=5r+1,l=r$ on one arm.  For $r\equiv3\pmod4$, this is the
E recipe of Section~\ref{templates:equal-r3}, with its finite
evaluations at $r=7,11$.  At $r=3$ the arm
$(15,6,12,7,14,4)$ has endpoint sum $19$.  For
$r\equiv1\pmod4$, $r\ge13$, use the alternative recipe of
Section~\ref{templates:alt}; its contracted component beginning with
$0,(5r+1)$ has length $2r$ and ends at $r$.
The same finite template is evaluated and checked at $r=9$.
The complete edge lists, original positions of their endpoints,
color-multiset cancellations, and finite weighted-component checks
are part of Section~\ref{templates:equal-even}; hence this endpoint
assertion is more than a degree or color calculation.  Apply the
operation just proved.

When $a$ is divisible by four and $a\ge12$, an equal-even source
with arms $a/2$ has balanced alpha index.  Corollary~\ref{common:balanced-alpha-matching} gives $\Sp{a,a,a+1}$:
its unique-odd-arm perfect matching has contree $\Sp{a/2,a/2,a/2}$
by Corollary~\ref{common:unique-odd-contree}.  The remaining targets $a=8,10$ have these
literal words, with cuts $12,15$, respectively:
\[
\begin{array}{l}
(25,1,24,2,23,3,21,4);\\
(19,6,20,8,17,10,18,7,22);\\
(16,11,14,12,13,9,15,5),
\end{array}
\qquad
\begin{array}{l}
(31,1,30,2,29,3,28,4,26,5);\\
(23,7,24,11,22,8,18,13,21,9,27);\\
(20,14,17,15,16,12,19,10,25,6).
\end{array}
\]
Their full label, color, and separation inventories are checked by
subtraction.  For completeness, the finite $r=9$ alternative evaluation
gives the following $a=18$ target with cut $27$:
\[
\begin{split}
 &(54,1,53,2,52,3,51,4,42,14,50,5,49,6,48,7,47,8);\\
 &(37,25,29,26,28,27,32,24,30,23,34,13,43,12,44,11,45,10);\\
 &(55,9,38,18,35,17,39,16,40,15,41,22,31,21,36,20,33,19,46).
\end{split}
\]
The cases just listed exhaust every even $a\ge6$.
\end{proof}

\subsection{Lower-hole kernels for two even long arms}
\label{mc:lower-kernels}

\begin{definition}\label{mc:lower-kernel}
A lower-hole kernel $L(s,u,t,H)$ is an alternating word of $2t$
vertices, starting upper and ending lower, with
\[
\begin{array}{ll}
\text{lower labels:}&\interval{1}{t+1}\setminus\{H\},\\
\text{upper labels:}&\interval{t+2}{2t+1},\\
\text{colors:}&\interval{1}{2t}\setminus\{D\},
 \quad D=s+2t-u+3-H,\\
\text{endpoints:}&E=2t-s+1,\quad F=t+u-2s-H.
\end{array}
\]
Only the indicated inventory conditions are imposed on auxiliary
source kernels.  The next lemma states the additional hypotheses
required to complete one to a spider.
\end{definition}

\begin{lemma}[Lower-hole completion]\label{mc:lower-completion}
Let $s\ge1$, $t\ge s$, $u\ge t$, and $1\le H\le t+1$.
Suppose $L(s,u,t,H)$ exists and
\[
 X=u+t-H+1\in\interval{u}{2u-1},\qquad
 2(X-1)\ne3(u-1).
\]
Then $\Sp{2s+1,2t,2u}$ is root-zero alpha with cut $s+t+u$,
and its shortest arm is the standard maximum Walecki word.
\end{lemma}
\begin{proof}
Map kernel lowers by $x\mapsto x+s$ and uppers by
$x\mapsto x+s+2u-1$, producing arm $A$.
Take the ordinary odd alpha path on $\interval{1}{2u-1}$
from upper endpoint $X$ to upper endpoint
$Y=3u-1-X=2u-t+H-2$.  Translate it by $s+t+1$ and append
the lower leaf $s+H$, obtaining arm $B$.

The internal colors of $A$ are
$\interval{2u}{2u+2t-1}\setminus\{R\}$, where
$R=D+2u-1=s+u+2t-H+2$.  Its root color is $2u+2t$.
The old internal colors of $B$ are $\interval{1}{2u-2}$,
its appended edge has color $Y+t+1-H=2u-1$, and its root color
is $X+s+t+1=R$.

For labels, $A$ supplies
\[
 (\interval{s+1}{s+t+1}\setminus\{s+H\})
 \ \cup\ \interval{s+t+2u+1}{s+2t+2u}.
\]
The lowers of $B$ fill the omitted $s+H$ and
$\interval{s+t+2}{s+t+u}$; its uppers fill
$\interval{s+t+u+1}{s+t+2u}$.
Thus $A,B$ partition $\interval{s+1}{s+2t+2u}$.
Set $N=2s+2t+2u+1$ and add
$W=(N,1,N-1,2,\ldots,N-s)$.  It supplies the two outer label
intervals and the color interval $\interval{2t+2u+1}{N}$.
All labels and colors occur once, and every edge crosses $s+t+u$.
The required strict bound on every label off $W$ also follows.
\end{proof}

\begin{lemma}[Diagonal growth of lower-hole kernels]
\label{mc:lower-growth}
Let $L(s_0,u_0,t_0,H)$ exist, put $g=t_0-s_0\ge1$, and
choose $P\ge g$.  Then $L(s_0+P,u_0+P,t_0+P,H)$ exists.
\end{lemma}
\begin{proof}
Take an even-order alpha path on $\interval{1}{2P}$ from $P+g$
to $g$, translate it by $t_0+1$, and prepend it to the old kernel
after raising the old uppers by $2P$.
Its last label is $t_0+1+g=2t_0-s_0+1$, and the mapped old
first label is $2P$ larger, so the join has color $2P$.
The colors are the new $\interval{1}{2P-1}$, that join, and
the old palette shifted by $2P$; the new missing color is $D_0+2P$.
The lower labels consist of the old lower interval with its hole
and $\interval{t_0+2}{t_0+P+1}$; the uppers form the two
adjacent intervals
\[
 \interval{t_0+P+2}{t_0+2P+1},\qquad
 \interval{t_0+2P+2}{2t_0+2P+1}.
\]
The new first endpoint is $E_0+P$ and its last is $F_0$.
These are exactly the required endpoints for the asserted parameters.
\end{proof}

\subsection{The complete distinct central OEE cone}
\label{mc:oee-central-section}

\begin{proposition}\label{mc:oee-central}
For positive odd $a$ and positive even $b,c$ satisfying
$a<b<c\le2a$, $\Sp{a,b,c}$ is root-zero alpha with its
shortest arm the standard maximum Walecki word.
\end{proposition}
\begin{proof}
Write $(a,b,c)=(2s+1,2t,2u)$, where
$1\le s<t<u\le2s+1$.
We first construct all points off the plane
\begin{equation}\label{mc:oee-plane}
 3u=2s+2t+3.
\end{equation}
Put $S=u-t-1$, $T=u-s-1$, $P=s+t-u+1$.
Then $0\le S<T$, $P\ge T-S\ge1$, $T+P=t$.
Take an even-order alpha path on $\interval{1}{2T}$
from $2T-S$ to $T-S$ and raise its uppers by $1+2P$.
Before it put an even-order alpha path on $\interval{1}{2P}$
from $P+T-S$ to $T-S$, translated by $T+1$.
The two endpoint requests are valid.  The join has color $2P$,
so the new word has missing color $D=2P+1$, lower hole
$H=T+1=u-s$, and endpoints $2t-s+1,t-s$.
Its full label classes are those of $L(s,u,t,H)$, since the
prefix supplies the intervening adjacent lower and upper blocks.

The completion request is $X=s+t+1$ on order $2u-1$.
Its endpoints $X$ and $3u-s-t-2$ lie in
$\interval{u}{2u-1}$.  The only forbidden upper request is
$2(X-1)=3(u-1)$, precisely \eqref{mc:oee-plane}.
Thus Lemma~\ref{mc:lower-completion} proves every off-plane point.

On the plane, first suppose $d=u-t\ge3$.  Set
\[
 R=s+t-u+2,\quad W=u-s-2,\quad A=t-s,\quad B=A+1.
\]
Here $R+W=t$, $1\le A\le R$ because $u\le2s+2$, and
$1\le B\le W$ because $d\ge3$.
Translate an even alpha path on $\interval{1}{2R}$ from
$R+A$ to $A$ by $W+1$.  Follow it with an even alpha path on
$\interval{1}{2W}$ from $W+B$ to $B$, keeping its lowers and
raising its uppers by $2R+1$.
The two lower intervals are $\interval{1}{W}$ and
$\interval{W+2}{t+1}$, so $H=W+1=u-s-1$.
The upper intervals adjoin to $\interval{t+2}{2t+1}$.
The internal color intervals are $\interval{1}{2R-1}$ and
$\interval{2R+2}{2t}$, while their join has color $2R+1$.
The missing color is $2R$, and the endpoints are
$2t-s+1$ and $t-s+1$, as required by $L(s,u,t,H)$.
The completion has $X=s+t+2$, $Y=3u-s-t-3$.
Both are upper; on \eqref{mc:oee-plane}, the forbidden equality
fails by two, since $2(X-1)=3u-1\ne3u-3$.
This fills every plane point with $d\ge3$.

For the remaining plane points put $j=s-(3d-2)$.  Then
\[
 s=3d-2+j,\quad t=3d-1+2j,\quad u=4d-1+2j,
 \qquad d\ge1,\ j\ge0.
\]
When $j\ge1$, $j\ne d+2,2d-2$, set
$T=d+j$, $P=2d-2+j$.
Take an odd alpha path on $\interval{1}{2T+1}$ with larger
class lower, from $d+2$ to $j$, and raise its uppers by $2P+2$.
These endpoints sum to $T+2$; its sole forbidden request is
$j=d+2$.  Before it put an odd alpha path on
$\interval{1}{2P+1}$ with larger class upper, from
$2d-1+2j$ to $4d-3+j$, translated by $T+2$.
Its endpoints sum to $3P+2$; its sole forbidden request is
$j=2d-2$.
The combined lower intervals are $\interval{1}{T+1}$ and
$\interval{T+3}{t+1}$, hence $H=T+2=u-s+1$.
The upper intervals adjoin to $\interval{t+2}{2t+1}$.
The two internal palettes are $\interval{1}{2P}$ and
$\interval{2P+3}{2t}$, and the join has color $2P+1$.
Thus the missing color is $2P+2$, with the required endpoints
$2t-s+1,j$.  The odd completion endpoints are
$6d-3+3j$ and $6d-1+3j$, both upper and distinct; their first
normalized endpoint differs by one from the forbidden midpoint.
There is no further exception.

For $d=1$ this leaves $j=0,3$; for $d=2$ it leaves $j=0,2,4$.
Four have the following literal lower-hole kernels:
\begin{center}\small
\begin{tabular}{ccl}
$(s,t,u)$&$H$&word\\\midrule
$(1,2,3)$&1&$(4,3,5,2)$\\
$(4,8,9)$&3&$(13,8,10,9,12,4,14,7,11,5,17,1,16,2,15,6)$\\
$(6,9,11)$&3&$(13,10,11,9,14,4,16,7,15,8,12,6,17,2,18,1,19,5)$\\
$(8,13,15)$&3&$(19,11,18,12,17,13,16,14,15,6,24,1,27,2,26,4,25,5,$\\
&&$22,8,23,7,20,10,21,9)$.
\end{tabular}\end{center}
Their completion orders and first labels are respectively
$(5,5),(17,15),(21,18),(29,26)$, all allowed.
The last point $(d,j)=(2,0)$ gives $(a,b,c)=(9,10,14)$;
its literal arm words, with cut $16$, are
\[
\begin{split}
 &(33,1,32,2,31,3,30,4,29);\\
 &(23,15,18,16,17,13,20,14,19,7);\\
 &(24,6,28,8,27,10,25,9,22,11,21,12,26,5).
\end{split}
\]
Full subtraction verifies every finite word above.  These branches
exhaust the plane, and all completions explicitly adjoin the required
Walecki arm.  Together with the off-plane construction they prove
the proposition, including its smallest instance $(3,4,6)$.
\end{proof}

\begin{corollary}\label{mc:eee-central}
Every all-even triple with $a<b<c\le2a+2$ is root-zero alpha.
\end{corollary}
\begin{proof}
The source $(a+1,b,c)$ satisfies the preceding proposition.
Shorten its prescribed Walecki arm by Lemma~\ref{mc:shorten}.
This includes $(2,4,6)$.
\end{proof}

\subsection{Repeated OEE triples}
\label{mc:oee-repeated-section}

\begin{lemma}[The adjacent repeated-even family]\label{mc:oee-near}
For even $b\ge6$, $\Sp{b-1,b,b}$ is root-zero alpha.
The shortest arm in this lemma is a modified Walecki word; the
standard-Walecki shortening lemma is not asserted for it.
\end{lemma}
\begin{proof}
Write $b=2u$, $u\ge3$, $N=6u-1$, $K=3u-1$.
If $u\ne5$, choose an ordinary alpha path on
$\interval{0}{2u-2}$ starting at $u+1$.  Its larger class is
upper and its other endpoint is $2u-4$.  Its only possible
forbidden endpoint request occurs at $u=5$: writing $u=2q+1$,
the upper forbidden label $3q$ equals $u+1=2q+2$ exactly at
$q=2$.  Add two to all labels and append the lower label $-1$.
The resulting word $Q$ has lowers
$\{-1\}\cup\interval{2}{u}$, uppers
$\interval{u+1}{2u}$, colors $\interval{1}{2u-1}$, first
label $u+3$, and last label $-1$.

Use three arms as follows.  First take the maximum Walecki word of
length $2u-1$, but raise its last lower label from $u-1$ to $u$:
\[
 W=(N,1,N-1,2,\ldots,5u+1,u,5u).
\]
Its lowers are $\interval{1}{u-2}\cup\{u\}$, its uppers
are $\interval{5u}{N}$, and its colors are
$\interval{4u+3}{N}\cup\{4u,4u+1\}$.
Map the lowers of $Q$ by $x\mapsto x+u$ and its uppers by
$x\mapsto x+3u-1$, giving an even arm $B$.
Its lowers are $\{u-1\}\cup\interval{u+2}{2u}$ and its
uppers are $\interval{4u}{5u-1}$.  It has internal colors
$\interval{2u}{4u-2}$ and root color $4u+2$.
For the other even arm $C$ use the pairs
\[
 (4u-1-j,2u+1+j)\quad(0\le j\le u-2),
 \qquad\text{followed by }(3u,u+1).
\]
Its lowers are $\interval{2u+1}{3u-1}\cup\{u+1\}$,
its uppers $\interval{3u}{4u-1}$, its internal colors
$\interval{1}{2u-1}$, and its root color $4u-1$.

These three lower sets partition $\interval{1}{K}$ and their
upper intervals partition $\interval{K+1}{N}$.  Their colors,
in order, partition
\[
 \interval{1}{2u-1},\quad\interval{2u}{4u-2},\quad
 \{4u-1,4u,4u+1,4u+2\},\quad\interval{4u+3}{N}.
\]
All edges cross $K$ and the root is zero.
At $u=5$ use the following literal words, with cut $14$:
\[
\begin{split}
 &(29,1,28,2,27,3,26,5,25);\\
 &(22,7,24,8,21,9,19,11,20,6);\\
 &(18,13,17,10,16,14,15,12,23,4).
\end{split}
\]
Their full inventories verify the omitted endpoint case directly.
\end{proof}

\begin{lemma}[Repeated central range]\label{mc:oee-repeated-central}
For $s\ge3$ and $2\le k\le s$, there is a root-zero alpha
labeling of $\Sp{2s+1,2s+2k,2s+2k}$ with standard Walecki
shortest arm.
\end{lemma}
\begin{proof}
For each $k\ge3$, the word
\[
 (2k+1-j,j+1)\quad(0\le j\le k-3),
 \qquad (k+3,k,k+2,k+1)
\]
is $L(0,k,k,k-1)$.  Its lower labels are
$\interval{1}{k-2}\cup\{k,k+1\}$; its uppers form
$\interval{k+2}{2k+1}$.  Its descending alternating initial
colors, the join to the final block, and that block itself give
\[
 \interval{6}{2k},\quad\{5\},\quad\{3,2,1\},
\]
respectively, with the obvious empty ranges at $k=3$.
Thus the missing color is four and its endpoints are $2k+1,k+1$.
Grow it by $P=s\ge k$ using Lemma~\ref{mc:lower-growth}.
The result has $t=u=s+k$, $H=k-1$, $D=2s+4$.
The completion request is
\[
 X=2s+k+2,\qquad Y=s+2k-3
\]
on order $2u-1$.  Both are in the upper interval because $k\ge3$;
the forbidden equality would require $k=s+5$, impossible when
$k\le s$.  Lemma~\ref{mc:lower-completion} applies.

For $k=2$ start instead with
$L(1,3,3,3)=(6,4,5,2,7,1)$ and grow by
$P=s-1\ge2$.  The resulting completion request on order
$2s+3$ has endpoints $2s+2,s+3$; its forbidden equality would
require $s=1$.  It is therefore allowed for every $s\ge3$.
\end{proof}

\begin{lemma}[The final repeated boundary]\label{mc:oee-repeated-boundary}
For $s\ge3$, $\Sp{2s+1,4s+2,4s+2}$ is root-zero alpha
with standard Walecki shortest arm.
\end{lemma}
\begin{proof}
Set $t=u=2s+1$.  Translate a standard maximum-Walecki odd path on
$\interval{1}{2s+1}$ by $s+2$.  It starts at $3s+3$, ends
at $2s+3$, has lower labels $\interval{s+3}{2s+2}$ and
upper labels $\interval{2s+3}{3s+3}$.
Follow it by an odd alpha path on $\interval{1}{2s+1}$ with
larger class lower, starting at two and ending at $s$, with its
uppers raised by $2s+2$.  Use the explicit endpoint-one word
of Section~\ref{templates:path-words}, shifted from zero-based
labels by one.  Its last three labels before the upper shift are
\[
 (s+1,s+2,s).
\]
For odd $s$ this follows from the final base $(1,2,0)$ shifted by
$s-1$ and then by one; for even $s$ it follows from the final
base tail $(4,5,3)$ shifted by $s-4$ and then by one.
These two cases apply for $s\ge3$ and $s\ge4$, respectively.
Its shifted internal colors are $\interval{2s+3}{4s+2}$.
The join has color $2s+1$; hence the combined kernel has hole
$H=s+2$, missing color $2s+2$, and endpoints $3s+3,s$.

Its last edge is $(3s+4,s)$, with color $2s+4$.
Change only its leaf $s$ to $s+2$.  This fills the old label hole
and opens the new hole $H'=s$.  The edge color becomes $2s+2$,
so the new missing color is $D'=2s+4$.
The new endpoint $s+2$ is exactly $t+u-2s-H'$.
Thus it is $L(s,2s+1,2s+1,s)$.

Its completion on order $4s+1$ has upper endpoints
$X=3s+3,Y=3s-1$.  The forbidden equality compares
$6s+4$ with $6s$, so cannot hold.  Apply
Lemma~\ref{mc:lower-completion} to obtain the claimed spider.
\end{proof}

\begin{lemma}[Two odd paths and reflection]\label{mc:two-odd-kernel}
Let $r,w\ge1$ and $\epsilon\in\{0,1\}$, and put
\[
 t=r+w+1,\quad u=s+w+1+\epsilon,\quad H=w+2.
\]
Define $X_A=2r+w-s+1$ and $X_B=s-r+2-\epsilon$.
If
\[
\begin{gathered}
 r+1\le X_A\le2r+1,\qquad 1\le X_B\le w+1,\\
 2(X_A-1)\ne3r,\qquad2(X_B-1)\ne w,
\end{gathered}
\]
then $L(s,u,t,H)$ exists.  Moreover, the class permutation
\[
 x\longmapsto t+2-x\ (x\le t+1),\qquad
 x\longmapsto3t+3-x\ (x\ge t+2)
\]
maps it to $L(t-s-1,2t-u+2,t,t+2-H)$.
\end{lemma}
\begin{proof}
Take an ordinary odd alpha path on $\interval{1}{2r+1}$
with larger class upper, starting at $X_A$ and ending at
$r-w+s+1$, and translate it by $w+2$.
Follow it by an odd alpha path on $\interval{1}{2w+1}$
with larger class lower, starting at $X_B$ and ending at
$w-s+r+\epsilon$; keep its lowers and raise its uppers by
$2r+2$.  The latter request is the complement of the ordinary
upper request, so its forbidden equality is exactly the one displayed.
The two paths give lower intervals $\interval{w+3}{t+1}$ and
$\interval{1}{w+1}$, and adjacent upper intervals
$\interval{t+2}{t+r+2}$ and
$\interval{t+r+3}{2t+1}$.
Their internal color intervals are $\interval{1}{2r}$ and
$\interval{2r+3}{2t}$, and the join has color $2r+1+\epsilon$.
Thus the missing color is $2r+2-\epsilon=D$, and direct substitution
gives the required endpoints.  The class permutation is a bijection
of the two label classes.  It changes each color $d$ to $2t+1-d$;
substitution in the endpoints, hole, and missing color gives exactly
the asserted new parameters.  It does not reverse the word.
\end{proof}

\begin{lemma}[The complete large repeated range]
\label{mc:oee-repeated-large}
For $S\ge3$, $T\ge2S+2$, $\Sp{2S+1,2T,2T}$ is root-zero
alpha with standard Walecki shortest arm.
\end{lemma}
\begin{proof}
Apply Lemma~\ref{mc:two-odd-kernel} with source parameters
\[
 s=T-S-1,\quad w=S+2-\epsilon,\quad r=T-w-1,
 \quad u=T+2.
\]
Its class reflection gives the target lower-hole parameters
$(S,T,T,H)$ with $H=r+1$.
The two source requests and the target completion are all valid if
\begin{equation}\label{mc:reflected-domain}
 T\ge2S+3-\epsilon,\quad T\ne3S+3-\epsilon,\quad
 S\ne4-3\epsilon,\quad T\ne2S+7-2\epsilon.
\end{equation}
For completeness, their first endpoints reduce to
$X_A=T-2+\epsilon$, $X_B=4-2\epsilon$, and
$X=T+S+3-\epsilon$.  The other completion endpoint is
$Y=2T-S-4+\epsilon$.
The three exclusions in \eqref{mc:reflected-domain} are respectively
the forbidden requests for the first source path, second source path,
and completion.  Thus none has been silently omitted.

For $\epsilon=1$, all $T\ge2S+2$ are covered except
$T=3S+2$ and $T=2S+5$.  For $T=3S+2$, the $\epsilon=0$
conditions hold unless $S=4$ or $S=5$.  For $T=2S+5$, they
hold unless $S=4$ (the other equality would force $S=2$).
Consequently exactly the following three finite kernels suffice:
\begin{center}\small
\begin{tabular}{ccl}
$(S,T)$&$H$&word\\\midrule
$(4,13)$&7&$(23,5,24,4,25,3,26,2,27,1,18,12,17,13,16,14,$\\
&&$15,8,22,6,21,9,20,10,19,11)$\\
$(4,14)$&9&$(25,7,24,5,26,6,22,8,23,1,29,2,28,3,27,4,$\\
&&$17,14,16,15,19,12,18,13,21,10,20,11)$\\
$(5,17)$&11&$(30,3,35,1,34,4,26,10,27,9,28,8,29,6,32,7,31,$\\
&&$2,33,5,20,18,19,16,21,17,23,14,22,15,25,12,24,13)$.
\end{tabular}\end{center}
These are $L(S,T,T,H)$ by direct complete inventory.
Their missing colors are $13,12,14$, respectively, and their
completion requests are $(2T-1,X)=(25,20),(27,20),(33,24)$.
All three requests are allowed.  Lemma~\ref{mc:lower-completion}
therefore fills these finite cases, proving the entire range.
\end{proof}

\begin{proposition}\label{mc:oee-repeated}
For every odd $a\ge7$ and even $b>a$, $\Sp{a,b,b}$ is
root-zero alpha.  Except at $b=a+1$, the constructions have the
standard maximum-Walecki shortest arm.
\end{proposition}
\begin{proof}
Write $a=2s+1,b=2u$, $s\ge3$.  If $u=s+1$ use
Lemma~\ref{mc:oee-near}.  If $s+2\le u\le2s$ use
Lemma~\ref{mc:oee-repeated-central}.  At $u=2s+1$ use
Lemma~\ref{mc:oee-repeated-boundary}.  For $u\ge2s+2$ use
Lemma~\ref{mc:oee-repeated-large}.  These disjoint ranges exhaust
all even $b>a$ and have the stated pinning properties.
\end{proof}

\begin{corollary}\label{mc:eee-repeated}
For even $a\ge6$ and even $b\ge a+4$, $\Sp{a,b,b}$ is
root-zero alpha.
\end{corollary}
\begin{proof}
Use Proposition~\ref{mc:oee-repeated} on $(a+1,b,b)$.
Here $b\ne(a+1)+1$, so the source has a standard Walecki
shortest arm and Lemma~\ref{mc:shorten} applies.  The modified
near-family word is not used in this deduction.
\end{proof}

\subsection{The complete EOO classification}
\label{construction:all-eoo-alpha3}
\label{mc:all-eoo-section}

\begin{theorem}\label{mc:all-eoo-large}
If $a\ge6$ is even and $a<b\le c$ are odd, then
$\Sp{a,b,c}$ is root-zero alpha.
\end{theorem}
\begin{proof}
Write $a=2s,b=2t+1,c=2u+1$, where $s\ge3$ and
$s\le t\le u$.  Every odd source shortened below is obtained
by a construction explicitly appending the maximum-Walecki arm of
length $2s+1$; an unpinned odd existence theorem is insufficient.
There are four exhaustive regions.

If $c\le2a+1$, Proposition~\ref{mc:eoo-central} applies.

If $c\ge a+b$, take the root-zero alpha pair $\Sp{a,b}$,
with cut $s+t$, and apply Lemma~\ref{common:alpha-extension}.
The only possible exception is
\begin{equation}\label{mc:eoo-far-resonance}
 c=4(s+t)+1.
\end{equation}
For it use the explicitly pinned repeated odd source
$\Sp{2s+1,2t+1,2t+1}$ from
Lemma~\ref{ok:repeated-kernels}.  It has a long upper leaf at
distance $e\le s+t+1$ above its cut $s+2t$.
Here is the full leaf-interface check.  At $t=s$, the equal-odd
formula has long leaves $5s+1,5s+2$ for $s\ge4$; at $s=3$
its literal equal-seven labeling has smaller long leaf $16$ and
cut nine.  At $t>s$, the generic repeated parameters give
\[
 e_A=t+2s-V+3,\qquad e_B=t+V.
\]
For $V=s+1$ choose $e_B=s+t+1$; for $V=s$ choose $e_B=s+t$;
for $V=s+3$ choose $e_A=s+t$.  These are exactly the three
parameter choices in the complete repeated construction.
Shorten its Walecki arm; all its other labels and these distances
are unchanged.  Append $2z$ vertices to the selected long leaf,
where $z=2s+t\ge e$, using
Lemma~\ref{common:upper-leaf-extension}.  The new length is
$2t+1+2(2s+t)=4(s+t)+1$, as required.

It remains that $c>2a+1$ and $c<a+b$.
Suppose first $b\ge2a+1$.  The explicit pinned families in
Lemma~\ref{ok:reverse-boundary} and Section~\ref{ok:short-wedges}
give this complete selection.  For $b=2a+1$, the middle half
is $2s$ and the largest half is $2s\le u\le3s-1$.
Its repeated endpoint uses the twice-ray kernel; the next point
$u=2s+1$ uses the class-reflected gap-two kernel, with its two
long carriers interchanged;
all $u\ge2s+2$ use the reversed wedge with its minus choice.
For $b=2a+3$, the middle half is $2s+1$ and the largest half
is at most $3s$; the plus-choice reversed wedge applies throughout.
For $b\ge2a+5$, put $U=(b-1)/2$, $\delta=(c-b)/2$.
Then $U\ge2s+2$ and $0\le\delta\le s-1$, exactly the
complete unequal-short-wedge domain, including its repaired
$\delta=s-3$ plane.  The indicated formulas, carrier orders,
and all their path bounds are proved in the cited sections.
Each completion pins the same Walecki short arm; shorten it.

Finally suppose $b<2a+1$ in that remaining region.  Then
\[
 u\ge2s+1,\quad t\le2s-1,\quad u\le s+t-1.
\]
Set $k=u-2s$, $h=s+t-u$.  Thus
\[
 k,h\ge1,\quad t=s+k+h,\quad u=2s+k,\quad s\ge k+h+1.
\]
The explicit transition construction in
Lemma~\ref{ok:transition-kernels} completes with a prescribed
Walecki shortest arm whenever $h\ge2$, or $h=1$ and
$s\ge k+3$.  For $h\ge2$ its source is the elementary kernel
$K(0,0,h,h+1)$, followed by carrier prefix and asymmetric growth.
The sole excluded boundary is $h=1,s=k+2$, whose source arms
are $(2s+1,4s-1,6s-3)$.
The controlled-edge induction in Lemma~\ref{ok:pinned-ap-ray}
constructs $K(s,3s-2,2s-1,2)$ for every $s\ge3$ and
completes it with the required pin.  This boundary also shortens.
The four regions exhaust the domain and prove the theorem.
\end{proof}

\begin{theorem}[Full EOO classification]\label{mc:all-eoo}
For positive integers $a<b\le c$, with $a$ even and $b,c$ odd,
$\Sp{a,b,c}$ is root-zero alpha if and only if
$(a,b,c)\ne(2,3,3)$.
\end{theorem}
\begin{proof}
The preceding theorem proves $a\ge6$.
For $a=2,4$ and $c\ge7$, use the complete minimum-two and
minimum-four constructions of Section~\ref{app:small-alpha}.
For $c<7$ the three positive cases have these literal arm words:
\[
\begin{array}{ll}
\Sp{2,3,5}:&(10,1);\ (6,2,9);\ (8,3,5,4,7),\\
\Sp{2,5,5}:&(12,1);\ (10,4,6,5,8);\ (9,2,7,3,11),\\
\Sp{4,5,5}:&(14,1,13,2);\ (10,3,8,4,12);\ (9,6,7,5,11).
\end{array}
\]
The full inventories follow by subtraction, with respective cuts
$4,5,6$ from the forced index formula.

For $\Sp{2,3,3}$ the cut would be three.
Every upper vertex is a neighbor of the root or a leaf, so none has
two nonroot lower neighbors.  Color one forces $(3,4)$.
Color two cannot be $(2,4)$ by this property, so it forces $(3,5)$.
Lower vertex three now has its full degree two occupied.
Color three can only be $(1,4)$, $(2,5)$, or $(3,6)$.
The first two violate the upper-vertex property, and the last gives
lower vertex three a third neighbor.  Each is impossible.
This excludes a root-zero alpha labeling of this profile and proves
the claimed classification.
\end{proof}

\subsection{The ten inputs and the five-to-four inputs}
\label{mc:input-summary}

The following table records exactly the constructions justifying
the seven-arm predicates.  The three lengths are sorted and at
least six.  Letters indicate their parities, and the forced cut is
always $\floor{a/2}+\floor{b/2}+\floor{c/2}$.
\begin{center}\small
\begin{tabular}{p{.17\linewidth}p{.42\linewidth}p{.32\linewidth}}
Input&Condition&Proof\\\midrule
Equal&$a=b=c$&Equal-triple formulas and equal-even repairs\\
Long&$c>2k$, $c\ne4k+1$, $k=\floor{a/2}+\floor{b/2}$
 &Lemma~\ref{common:alpha-extension}\\
OOO central&$c\le2a-1$&Proposition~\ref{ok:central-cone}\\
EOO central&$c\le2a+1$&Proposition~\ref{mc:eoo-central}\\
EEO central&$a<b$, $c\le2a-1$&Corollary~\ref{mc:eeo-central}\\
EEO near&$a=b$, $c=a+1$&Proposition~\ref{mc:eeo-near}\\
OEE central&$a<b<c\le2a$&Proposition~\ref{mc:oee-central}\\
OEE repeated&$a<b=c$&Proposition~\ref{mc:oee-repeated}\\
EEE central&$a<b<c\le2a+2$&Corollary~\ref{mc:eee-central}\\
EEE repeated&$a<b=c$, $b\ge a+4$&Corollary~\ref{mc:eee-repeated}.
\end{tabular}\end{center}
The equal-odd formula is in Proposition~\ref{ok:equal-odd-three};
equal-even repairs and their exact small bases are in
Section~\ref{templates:equal-even} and Section~\ref{app:finite-bases}.
None uses a seven-arm selection theorem.

The manual five-to-four proof additionally uses
Proposition~\ref{mc:oeo-central} and the EOO near family
$\Sp{a,a+1,a+1}$ from Proposition~\ref{mc:eoo-repeated}.
It can use the smaller pre-repair EOO central domain, which is
contained in Proposition~\ref{mc:eoo-central}.
Deriving the six-to-four bridge from that manual lemma therefore
introduces no circularity: the construction layer precedes every
seven-arm selection and unit-extension argument.

Every infinite assertion here reduces to prescribed-endpoint paths,
explicit interval operations, a complete induction, or the
finite-template verification theorem.  Printed finite words and
finite-template tables form part of those proofs.  Independent graph
and integer audits provide additional verification; bounded samples
are not used to infer an infinite family.
